\documentclass[11pt,letter]{article}

\usepackage[margin=2.5cm]{geometry}
\usepackage{amsmath}
\usepackage{amssymb}
\usepackage{bm}
\usepackage{booktabs}
\usepackage{array}
\usepackage{ragged2e}
\usepackage{seqsplit}
\usepackage{enumitem}
\usepackage{xcolor}
\usepackage[hidelinks]{hyperref}
\usepackage{setspace}
\setstretch{1.15}

\setlength{\emergencystretch}{3em}
\newcolumntype{L}[1]{>{\RaggedRight\arraybackslash}p{#1}}

%----------------------------------------------------------------------
%  Text helpers
%----------------------------------------------------------------------
\newcommand{\file}[1]{\texttt{\seqsplit{#1}}}
\newcommand{\code}[1]{\texttt{#1}}
\newcommand{\SCALAR}{\textbf{scalar}}
\newcommand{\ToClaude}[1]{[\textbf{CLAUDE}: {#1}]}
\newcommand{\status}[1]{\par\noindent\textit{Status:} #1\par}

%----------------------------------------------------------------------
%  Notation (Section 2).  Arrays bold, tables blackboard bold, scalars
%  plain.  Accents: \bar = mean over a region, \widetilde = smoothed,
%  \widehat = fitted or estimated.
%----------------------------------------------------------------------
% images on the pixel grid, one channel (bold uppercase)
\newcommand{\mE}{\mathbf{E}}   % lit layer as read, DN
\newcommand{\mF}{\mathbf{F}}   % dark-corrected signal, DN
\newcommand{\mP}{\mathbf{P}}   % flat-field polynomial shape
\newcommand{\mG}{\mathbf{G}}   % gray-chip level map, DN
\newcommand{\mR}{\mathbf{R}}   % post-image reflectance vs gray chips
\newcommand{\mS}{\mathbf{S}}   % pre-image reflectance, lined up
\newcommand{\mT}{\mathbf{T}}   % shading-corrected reflectance
\newcommand{\mA}{\mathbf{A}}   % post-image absorbance
\newcommand{\mB}{\mathbf{B}}   % pre-image absorbance
\newcommand{\mC}{\mathbf{C}}   % stain absorbance (chain output)
\newcommand{\mU}{\mathbf{U}}   % yellow amount image
\newcommand{\mV}{\mathbf{V}}   % brown amount image
\newcommand{\tB}{\widetilde{\mathbf{B}}}   % B smoothed
\newcommand{\hB}{\widehat{\mathbf{B}}}     % estimated shading of the pad
\newcommand{\hS}{\widehat{\mathbf{S}}}     % per-pixel white
% pixel x channel tables (blackboard bold)
\newcommand{\tA}{\mathbb{A}}
\newcommand{\tC}{\mathbb{C}}
% 1-D arrays (bold lowercase)
\newcommand{\vc}{\mathbf{c}}           % one pixel's spectrum
\newcommand{\vone}{\mathbf{1}}
\newcommand{\epsu}{\bm{\varepsilon}_u} % yellow spectrum
\newcommand{\epsv}{\bm{\varepsilon}_v} % brown spectrum
% regions: a text identifier set under the symbol
\newcommand{\reg}[2]{\underset{\text{#1}}{#2}}
\newcommand{\avg}[2]{\overline{\reg{#1}{#2}}}
% shading methods
\newcommand{\MA}{\mathcal{A}}
\newcommand{\MB}{\mathcal{B}}
\newcommand{\MC}{\mathcal{C}}
\newcommand{\MD}{\mathcal{D}}

\title{The analysis chain, from photo to arsenic\\
\large The steps that make the number, one notation for all of them,\\
and what is inspection or research}
\author{Colorimeter project --- pad analysis}
\date{Draft for Bill's review, 2026-10-03.  Branch \file{preimage-capture},
commit \file{4962604} plus uncommitted work.  No chain code is changed by
this document.}

\begin{document}
\maketitle

\begin{abstract}
\noindent
The chain that turns a photo of an indicator strip into an arsenic
estimate has a few simple steps.  The research behind them has left
long code and mixed names.  This document is the spine for the rest of the
write-ups.  It states the steps that make the number, one short section
each, in a single notation where an array can't be mistaken for a single
number.  It lists every other piece of code as inspection, research or
history.  And it lists what is still to be settled.  The worked numbers
follow two pixels of C1032, channel royal\_blue.B, through every step.
\end{abstract}

\tableofcontents
\newpage

%======================================================================
\section{What this document is for}

Three things, in this order:
\begin{enumerate}
  \item \textbf{The chain} (Section~\ref{sec:chain}): the steps needed to get
    from the stored photo to an arsenic concentration, and nothing else.
    Each step has a short description in words, one equation, a worked
    number, the code that does it today, and every scalar it uses.
  \item \textbf{The notation} (Section~\ref{sec:notation}): one symbol per
    quantity, used in every later write-up.  Tangential studies may reuse
    letters; the chain may not.
  \item \textbf{What is not the chain} (Section~\ref{sec:inventory}) and
    \textbf{what is still open} (Section~\ref{sec:open}).
\end{enumerate}

\noindent The exposure geometry (how the strip meets the arsine) is not
final.  Early strips were exposed over the whole pad, recent ones in a dot,
and the dot may change.  Each geometry shows a different fraction of the
total exposure, so the calibration found now may not carry over.  The chain
below avoids depending on the dot's size or shape until the spatial step.

%======================================================================
\section{Light adds; factors multiply}
\label{sec:rules}

Two kinds of operation run through the chain: \emph{combining} light, and
\emph{removing} an influence on it.  A simple thought experiment shows which
arithmetic each needs.  The letters $n$, $a$, $b$, $d$ are used in this
section only.

\subsection{The thought experiment}

Follow the photons that end up in one pixel:
\begin{itemize}[nosep]
  \item the LEDs emit $n$ photons toward the pad (set by the exposure time
    and the lamp's brightness);
  \item a fraction $b$ of that illumination reaches this spot.  The pad is
    lit by four LEDs, one beyond each corner, and their light adds, so $b$
    is the sum of what the LEDs that can see the spot deliver.  $b = 1$ when
    the spot is fully lit; a \emph{shadow} is $b < 1$: light never added,
    because one or more LEDs are blocked from the spot;
  \item a fraction $a$ of the light that arrives is sent back into the
    sensor.  $a$ belongs to the surface: the paper, and any dye on it, which
    lowers $a$.  A shadow does not change $a$.
\end{itemize}
The photons captured are
\begin{equation}
  d = a\, b\, n .
\end{equation}
Now take the same spot in the two photos.  The strip and the LEDs are in
the same places, so the shadow is the same, $b$; the surface changed from
$a_\text{pre}$ to $a_\text{post}$; the exposures differ, $n_\text{pre}$
and $n_\text{post}$:
\begin{equation}
  d_\text{pre} = a_\text{pre}\, b\, n_\text{pre}, \qquad
  d_\text{post} = a_\text{post}\, b\, n_\text{post}.
\end{equation}

\noindent\textbf{The ratio} cancels the shadow, whatever the dye:
\begin{equation}
  \frac{d_\text{post}}{d_\text{pre}}
  = \frac{a_\text{post}}{a_\text{pre}} \cdot
    \frac{n_\text{post}}{n_\text{pre}} .
\end{equation}
$b$ is gone.  $n_\text{post}/n_\text{pre}$ is the exposure and lamp
difference, which the gray chips measure (step 3, and step 6, ``What the
gray-chip level does'').  What is left, $a_\text{post}/a_\text{pre}$, is the
stain.

\noindent\textbf{The difference} does not:
\begin{equation}
  d_\text{post} - d_\text{pre}
  = b\,\bigl(a_\text{post}\, n_\text{post} - a_\text{pre}\, n_\text{pre}\bigr).
\end{equation}
$b$ stays as a multiplier, so the shadow stays in the answer.

\noindent\textbf{Worked example.}  Equal exposures; the stain halves the
light ($a_\text{pre} = 1$, $a_\text{post} = 0.5$); in the shadow one of the
four LEDs is blocked ($b = 0.75$):

\medskip
\begin{center}
\begin{tabular}{@{}lrrrr@{}}
\toprule
 & $d_\text{pre}$ & $d_\text{post}$ & ratio & difference \\
\midrule
fully lit ($b = 1$)    & 1.00 & 0.50  & \textbf{0.50} & 0.50 \\
in shadow ($b = 0.75$) & 0.75 & 0.375 & \textbf{0.50} & 0.375 \\
\bottomrule
\end{tabular}
\end{center}

\noindent The ratio reads the same stain in and out of the shadow; the
difference of counts reads the shadowed spot as 25\,\% less stained.

\subsection{The two rules}

\textbf{Rule 1: combining light is an ordinary (arithmetic) average or sum
of counts or reflectance.}  Light adds: from different LEDs (that is what
$b$ is), from different places, and from different moments.  A pixel
records the light collected over its area, so a pixel half in shadow
records the average of $b$ over the pixel.  Any step that stands for ``the
light over a larger or a shifted area'' --- averaging neighboring pixels,
lining up the pre-image (step 4), smoothing --- is therefore an arithmetic
average of counts, or of reflectance, which is counts divided by a smooth
gray level.  Averaging absorbances instead amounts to a geometric mean of
the reflectances, which gives the darker pixels more weight than the light
they actually collected.

\textbf{Rule 2: removing an influence that enters as a factor, as the
shadow $b$, the exposure $n$ and the dye's $a$ do in $d = a\,b\,n$, is a
division in counts or reflectance, which is exactly a subtraction in
absorbance.}  Because the logarithm of a product is the sum of the
logarithms,
\begin{equation}
  -\log_{10} d = -\log_{10} a - \log_{10} b - \log_{10} n ,
\end{equation}
dividing by $b$ and subtracting $-\log_{10} b$ give identical numbers.  So
the shadow \emph{is} removed in reflectance and in absorbance at the same
time; the chain writes it as a subtraction from step 5 on only because it
works in absorbance there.  The same is true of the exposure (through the
gray chips) and of the stain.

\textbf{One limit.}  The thought experiment treats one spot.  A pixel
averages over its area (rule 1), so it records the average of $a\,b$, which
equals (average of $a$)$\times$(average of $b$) only if the dye or the shadow
is even across that pixel.  The two differ only where a shadow's edge and a
stain's edge fall in the same pixel.  On the current dot samples the dot
lies in the unshaded inner pad, so the ratio and the difference of counts
would give nearly the same answer there; the difference matters for
whole-pad exposure, or for a dot that reaches the rim.

\subsection{Why work in absorbance after step 5}

With every factor turned into a term, the stain becomes a sum: a shading
term, plus a yellow term, plus a brown term (dyes add approximately in
absorbance; step 8).  The yellow/brown split solves sums.  A shadow blocks
every LED color alike, so its term is the same in every channel, while each
dye's term changes from channel to channel; that is how the split tells
them apart.

\subsection{Where the domain matters: only in averages}

Division and subtraction do not care which domain they are written in;
averaging does.  The chain averages in two places.  Lining up (step 4)
averages reflectance, as rule 1 says.  Smoothing the pre-image (step 6) is
at present an average of absorbances.  On C1032 royal\_blue.B the
difference between the two is under 0.0002 more than 4 pixels in from the
pad's edge ($\sigma = 1$ px), but up to 0.008 (median 0.0004--0.0008)
within 4 pixels of the edge, where the shadow changes quickly
(\file{chain\_example.smooth\_mean\_check}).  The smoothed pre-image stands
for the illumination over a larger area, so by rule 1 it should be an
average of reflectance; see step 6 and Section~\ref{sec:open}.

%======================================================================
\section{Notation}
\label{sec:notation}

\subsection{Conventions}

\begin{tabular}{@{}L{4.6cm}L{10.2cm}@{}}
\toprule
Style & Meaning \\
\midrule
$\mA$ (bold uppercase) & An \emph{image}: one value per pixel on the
  photo's pixel grid, for one channel.  130 rows by 183 columns for the
  current crop. \\
$A_{ij}$ (plain, subscripted) & One element of that image: the pixel in row
  $i$, column $j$.  A single number. \\
$\tA$ (blackboard bold) & A \emph{table}: one row per pixel, one column per
  channel ($N \times K$).  Same letter as the image it is made from. \\
$\vc_n$ (bold lowercase) & A list of numbers, e.g.\ one row of a table: one
  pixel's values over the $K$ channels. \\
$e$, $\lambda$ (plain) & A \SCALAR: one number.  Every scalar the chain
  applies to many pixels is flagged in bold in the text, because it assumes
  those pixels are alike. \\
$\reg{pad}{\mA}$ & The image restricted to a region; the region's name is
  written under the symbol. \\
$\avg{g1}{\mF}$ & Overline: the mean over that region.  A scalar. \\
$\tB$ & Tilde: smoothed.  An image. \\
$\hB$ & Hat: fitted or estimated.  An image. \\
$\MA, \MB, \MC, \MD$ & Calligraphic: the four shading-removal methods of
  \file{shading\_removal.tex}. \\
\bottomrule
\end{tabular}

\medskip
\noindent Arithmetic between images is done pixel by pixel.  So
``$\mR = \mF / \mG$'' means $R_{ij} = F_{ij} / G_{ij}$ at every pixel; no
matrix product is meant.  Equations below give the element form.

\subsection{Indices and regions}

\begin{tabular}{@{}L{2.4cm}L{12.4cm}@{}}
\toprule
$i$, $j$ & Row and column of a pixel. \\
$k$ & Channel: one LED and one Bayer layer, e.g.\ royal\_blue.B.  $K$
  channels in all. \\
$n$ & Pixel number in a list (a row of a table), $1 \dots N$. \\
$\ell$ & Strip (sample), $1 \dots L$. \\
\midrule
pad & The fitted pad rectangle, hot pixels left out. \\
inner, rim & Pad pixels more than, and within, $w$ pixels of the pad's edge. \\
dot, free & Inside the dot's edge; at least 6 pixels outside it (dye-free).
  Geometry not final. \\
g1, g2 & The left and right gray chips. \\
\bottomrule
\end{tabular}

\subsection{The symbols, by group}

Letters come in groups, so the group tells what kind of quantity it is:
digital number (DN) $E, F, G, P$; 
reflectance $R, S, T$; 
absorbance $A, B, C$; 
yellow and brown $U, V$.

\medskip
\begin{tabular}{@{}L{1.5cm}L{7.3cm}L{1.6cm}L{3.6cm}@{}}
\toprule
Symbol & Meaning & Kind & Code name today \\
\midrule
\multicolumn{4}{@{}l}{\emph{Counts (DN, the camera's raw unit, 0--1023)}} \\
$\mE$ & the lit layer as read from the .h5 file & image & \code{raw}, \code{lit} \\
$e$ & dark level: the camera's reading with no light & \SCALAR & \code{dark\_level} \\
$\mF$ & dark-corrected signal, $F_{ij} = E_{ij} - e$ & image & \code{sig} \\
$\mP$ & flat-field polynomial: the shape of the light on the stage, mean 1, measured once & image & \code{gain\_at(...)} \\
$\gamma_1$, $\gamma_2$ & light level at each gray chip: mean of $\mF$ over the chip / mean of $\mP$ there & \SCALAR{}s & \code{s1/g1}, \code{s2/g2} \\
$\gamma_j$ & the level blended between the chips, one per column & 1-D & \code{level} \\
$\mG$ & gray-chip level, $G_{ij} = P_{ij}\gamma_j$ & image & \code{div} \\
\midrule
\multicolumn{4}{@{}l}{\emph{Reflectance (no units; bigger = lighter)}} \\
$\mR$ & post-image, divided by the gray-chip level & image & \code{R} \\
$\mS$ & pre-image, the same steps, lined up on $\mR$'s grid & image & \code{R} (pre, registered) \\
$\hS$ & the per-pixel white: $10^{-\hB}$ & image & (new) \\
$\mT$ & shading-corrected reflectance, $T_{ij} = R_{ij} / \widehat S_{ij}$ & image & (new) \\
\midrule
\multicolumn{4}{@{}l}{\emph{Absorbance (no units; bigger = darker)}} \\
$\mA$ & post-image, $A_{ij} = -\log_{10} R_{ij}$ & image & (A without \code{R\_white}) \\
$\mB$ & pre-image, $B_{ij} = -\log_{10} S_{ij}$ & image & \code{A\_pre} \\
$\tB$ & $\mB$ smoothed, width $\sigma$ & image & \code{smooth(A\_pre)} \\
$\hB$ & estimated shading: $\lambda\mB + (1-\lambda)\tB$ & image & \code{corr} (method D) \\
$\mC$ & the stain: $C_{ij} = A_{ij} - \widehat B_{ij}$.  \textbf{The chain's
  photometric output.} & image & \code{A\_after} \\
$\tC$ & $\mC$ for all channels, pad pixels only & table & \code{PadCube.A} \\
\midrule
\multicolumn{4}{@{}l}{\emph{Yellow and brown (under review)}} \\
$\epsu$, $\epsv$ & the yellow and brown reference spectra, one value per channel & list & \code{eps\_y}, \code{eps\_b} \\
$\mU$, $\mV$ & yellow and brown amounts at each pixel & image & \code{y}, \code{b} \\
\midrule
\multicolumn{4}{@{}l}{\emph{Settings and calibration}} \\
$\lambda$ & weight on the unsmoothed pre-image & \SCALAR & \code{wt} \\
$\sigma$ & smoothing width, pixels & \SCALAR & \code{R3\_D\_PX} \\
$w$ & rim width, pixels & \SCALAR & \code{INNER\_PX} \\
$x_\ell$ & arsenic concentration of strip $\ell$'s standard, ppb & \SCALAR & \code{manual\_read} \\
\bottomrule
\end{tabular}

\medskip
\noindent Two pairs are worth remembering: $\mR \leftrightarrow \mA$ and
$\mS \leftrightarrow \mB$ (the same photo in reflectance and in
absorbance), and $\mT \leftrightarrow \mC$ (the corrected result,
$\mT = 10^{-\mC}$).

%======================================================================
\section{The chain}
\label{sec:chain}

Everything in this section is needed to get from the stored photo to an
arsenic concentration.  Steps 1--6 are done one channel at a time.  The
worked numbers are for C1032, royal\_blue.B, at two pixels, each followed
through both photos (post-image and pre-image):
\begin{itemize}[nosep]
  \item \textbf{a dye-free pixel}, row 55, column 78: 8.8 pixels in from the
    pad's edge and 7.3 pixels outside the dot;
  \item \textbf{a dot pixel}, row 66, column 91: near the dot's center.
\end{itemize}
They are printed by \file{chain\_example.worked\_example()}
(\file{desktop/padanalysis/chain\_example.py}).

\subsection*{Before step 1: where things are (geometry)}

Three things are found once per photo and used by every channel:
\begin{itemize}
  \item \textbf{Hot pixels}: positions whose dark reading is reliably high.
    They are left out, not repaired (\code{probe.Capture.hot\_mask}, from
    the \code{hotpixels} database table).  C1032: 8 hot pixels on the pad, leaving 2215 usable pad pixels.
  \item \textbf{The pad}: a rotated rectangle of fixed size, found from the
    strip's edges (\file{padrect.PadRect}).  It gives each pixel's depth
    into the pad, which defines the inner pad and the rim.
  \item \textbf{The dot}: its edge, from the post-image
    (\code{pca\_study.dot\_geometry}).  Needed only for the spatial step and
    for choosing the dye-free region in studies.  Geometry not final.
\end{itemize}

%----------------------------------------------------------------------
\subsection{Step 1: read the layer and choose the exposure}

\textbf{In words.}  The camera took each LED's photo at one or more
exposure times.  For each channel, use the brightest exposure that is not
clipped anywhere on the pad or the two gray chips and that carries more
than read noise.

\textbf{Equation.}  $\mE$ = the stored layer for (LED, Bayer layer,
chosen exposure).

\textbf{Worked numbers.}  Post-image: rung e1, 2193\,\textmu s.  Pre-image:
rung e0, 1000\,\textmu s.  In every table of this section the left column is
the dye-free pixel and the right column the dot pixel.  The pre-image's
values are taken where that post-image pixel came from (step 4), by the
same 4-pixel average the chain uses.

\medskip
\begin{center}
\begin{tabular}{@{}lrr@{}}
\toprule
 & dye-free (55, 78) & dot (66, 91) \\
\midrule
post-image $E$, DN & 782.00 & 266.00 \\
pre-image $E$, DN  & 396.99 & 394.06 \\
\bottomrule
\end{tabular}
\end{center}

\noindent The pre-image reads about half the post-image's counts at the
dye-free pixel because its exposure is less than half as long.

\textbf{Code today.}  \code{probe.Capture.layer}; the choice is
\code{analysis\_1.rung\_verdict}.  Clipped means $\ge 1000$\,DN
(\code{SATURATED\_DN}); read noise means a mean below 25\,DN
(\code{MIN\_SIGNAL\_DN}).

\textbf{Scalars.}  The two thresholds.  The pre-image and post-image may be
at different exposures; step 3 removes that.

\status{in use.  Open: blending several exposures into one (Bill to
review; Section~\ref{sec:open}).}

%----------------------------------------------------------------------
\subsection{Step 2: subtract the dark level}

\textbf{In words.}  Take away what the sensor reads with no light.

\textbf{Equation.}
\begin{equation}
  F_{ij} = E_{ij} - e
\end{equation}
$e$ is a \SCALAR: the average of one dark frame, one number per Bayer
layer per photo (\file{tex/dark\_signal\_characterization.tex}).

\textbf{Worked numbers.}  $e = 64.028$\,DN (post-image), $64.021$\,DN
(pre-image).

\medskip
\begin{center}
\begin{tabular}{@{}lrr@{}}
\toprule
 & dye-free & dot \\
\midrule
post-image $F$, DN & 717.97 & 201.97 \\
pre-image $F$, DN  & 332.96 & 330.04 \\
\bottomrule
\end{tabular}
\end{center}

\textbf{Code today.}  \code{probe.Capture.layer} returns $\mE - e$;
$e$ comes from \code{capture\_h5.dark\_level}.

\status{settled (dark-signal report).}

%----------------------------------------------------------------------
\subsection{Step 3: divide by the gray-chip level (reflectance)}

\textbf{In words.}  The light is not the same everywhere on the stage, and
the lamp's brightness and the exposure time change from photo to photo.
The flat-field polynomial $\mP$ gives the light's \emph{shape}, measured
once.  The two gray chips, read in this same photo, give its \emph{level}.
Dividing by the result, $\mG$, makes each pixel's brightness relative to the
gray reference at that place.  The lamp's brightness and the exposure time
scale $\mF$ and $\mG$ alike, so they drop out of the ratio.  $\mG$ is a
smooth surface (a low-order polynomial set by two chip means), so it adds
almost no pixel noise, but it removes none either: all of $\mF$'s pixel
noise and paper texture stay in $\mR$.

\textbf{Equations.}  First the light level at each chip, then a level for
each column, blended in a straight line from the left chip (center column
$j_1$) to the right chip (center column $j_2$):
\begin{align}
  \gamma_1 &= \frac{\avg{g1}{\mF}}{\avg{g1}{\mP}}, \qquad
  \gamma_2 = \frac{\avg{g2}{\mF}}{\avg{g2}{\mP}}
    && \text{(\SCALAR{}s, one pair per photo)} \\
  t_j &= \frac{j - j_1}{j_2 - j_1}, \qquad
  \gamma_j = (1 - t_j)\,\gamma_1 + t_j\,\gamma_2
    && \text{(one value per column)} \\
  G_{ij} &= P_{ij}\,\gamma_j, \qquad R_{ij} = F_{ij} / G_{ij}
\end{align}
$\gamma_1$ is ``how many counts per unit of the polynomial at the left
chip''.  The chip means of $\mP$ are the same in every photo (the polynomial
and the chip rectangles are fixed); the chip means of $\mF$ are new in each
photo, two for the post-image and two for the pre-image.  $R = 1$ means ``as
bright as the gray chips would be here''.  The pad reads about 2, because
it is about twice as bright as the chips.

\textbf{Worked numbers.}  Chip means of $\mP$: $\avg{g1}{\mP} = 1.0052$,
$\avg{g2}{\mP} = 0.9933$; $j_1 = 26.5$, $j_2 = 160.5$.

\medskip
\begin{center}
\begin{tabular}{@{}lrr@{}}
\toprule
 & left chip (g1) & right chip (g2) \\
\midrule
post-image $\avg{g}{\mF}$, DN & 331.62 & 331.39 \\
post-image $\gamma$, DN       & 329.92 & 333.61 \\
pre-image $\avg{g}{\mF}$, DN  & 149.46 & 149.47 \\
pre-image $\gamma$, DN        & 148.70 & 150.47 \\
\bottomrule
\end{tabular}
\qquad
\begin{tabular}{@{}lrr@{}}
\toprule
 & dye-free & dot \\
\midrule
$P_{ij}$          & 1.0160 & 1.0172 \\
$t_j$             & 0.3844 & 0.4814 \\
post $\gamma_j$, DN & 331.34 & 331.69 \\
post $G_{ij}$, DN   & 336.65 & 337.40 \\
post $R_{ij}$       & 2.1327 & 0.5986 \\
pre $\gamma_j$, DN  & 149.38 & 149.55 \\
pre $G$, DN         & 151.77 & 152.12 \\
pre $R$             & 2.1939 & 2.1696 \\
\bottomrule
\end{tabular}
\end{center}

\noindent Check, post-image dye-free pixel:
$\gamma_{78} = 0.6156 \times 329.92 + 0.3844 \times 333.61 = 331.34$;
$G = 1.0160 \times 331.34 = 336.65$; $R = 717.97 / 336.65 = 2.1327$.
The polynomial values $P$ for the pre-image are the same to four digits,
because the strip moved less than a pixel.

\textbf{Why two chips (a choice to revisit).}  Two chips allow the level
to differ from left to right, in the column direction only.  The choice can
be revisited once the final reference color chips are installed and
calibrated: a single gray value, averaged from the left and right chips, may
be enough; or gray could be measured at more than two places around the pad
and the polynomial adjusted with a higher-order correction.  For the stain
$\mC$, what matters is set out in step 6: on C1032 a single averaged chip
would change $\mC$ by at most 0.0003.

\textbf{Code today.}  \code{padphot.divisor\_map} ($\mG$),
\code{padphot.reflectance} ($\mR$); $\mP$ is
\file{whitecorr.LayerModel.gain\_at}.

\status{in use (\file{tex/colorimeter\_flatfield.tex}).  That report uses
older letters; see Section~\ref{sec:open}.}

%----------------------------------------------------------------------
\subsection{Step 4: the pre-image, the same way, lined up}

\textbf{In words.}  The pre-image (the same strip, photographed before
exposure to arsine) goes through steps 1--3 \emph{in its own frame}, with
its own gray chips, because the light belongs to the stage, not to the
strip.  Only then is it moved onto the post-image's grid, because the strip
was taken out and put back between the two photos.
\underline{Each lined-up value is a weighted average of the 4 nearest}
\underline{pre-image pixels.}

\textbf{Equation.}  $\mS$ = the pre-image's $\mR$, lined up on the
post-image's grid.

\textbf{Worked numbers.}  Larger is brighter; the dot pixel was as bright as
the dye-free one before exposure.

\medskip
\begin{center}
\begin{tabular}{@{}lrr@{}}
\toprule
 & dye-free & dot \\
\midrule
taken from pre-image position $(i, j)$ & (54.83, 78.00) & (65.83, 91.03) \\
$S_{ij}$ & 2.1939 & 2.1696 \\
\bottomrule
\end{tabular}
\end{center}

\textbf{Is it right to average reflectances?}  Yes.  A pixel collects
light over its area, and reflectance is light, so the light a pixel at an
in-between position would have collected is the weighted average of the
reflectances around it.  Averaging absorbances instead would be a geometric
mean of the reflectances, which gives less weight to the bright pixels.  The
difference only shows where neighbors differ a lot: on C1032 it is at most
0.0002 in absorbance more than 1 pixel in from the pad's edge, and up to
0.013 (median 0.0008) at the outermost pixel, where the edge shadow changes
by up to 30\,\% from one pixel to the next
(\file{chain\_example.lineup\_mean\_check}).

\textbf{The averaging lowers the pixel-to-pixel scatter.}  Measured on 18
pads and six channels, the lined-up pre-image keeps a median 0.61 of its
scatter (range 0.43--0.96), depending on where each pad landed between
pixels (\file{shading\_removal.tex}, Section 4).  So noise is measured on
the pre-image's own grid.

\textbf{Code today.}  \code{analysis\_1.get\_balanced\_cell\_data(...,
sec='pre')}; the move is \code{analysis\_1.pre\_to\_post} and
\code{\_register} (from the two fitted pad rectangles; no image
matching).

\status{in use.  Open: how to allow for the pad-to-pad difference in
averaging.  Bill (2026-10-04): compare (a) recording each pad's fractional
position and allowing for it when setting $\sigma$ and $\lambda$, with (b)
lining up by nearest pixel; decide in the rim and edge-smoothing step.}

%----------------------------------------------------------------------
\subsection{Step 5: absorbance}

\textbf{In words.}  Change to a scale where a shadow, and a dye, add rather
than multiply.  The minus sign makes darker give a bigger number: a
reflectance below 1 has a negative logarithm, and the minus turns it
positive.\footnote{The form comes from light passing through a dye: the
fraction that gets through is $10^{-\alpha}$, where $\alpha$ grows in
proportion to the amount of dye (the Beer--Lambert law), so
$\alpha = -\log_{10}(\text{fraction through})$ measures the amount of dye.  For light scattered back out of
paper this holds only roughly.  See \file{duotone\_guide.tex}, ``Why
absorbance, not reflectance?''.}

\textbf{Equations.}
\begin{equation}
  A_{ij} = -\log_{10} R_{ij}, \qquad B_{ij} = -\log_{10} S_{ij}
\end{equation}
On the pad these are negative, because $R$ and $S$ are measured against the
gray chips and the pad is brighter than the chips.  That is expected: only
their difference, in step 6, is the stain.

\textbf{Worked numbers.}

\medskip
\begin{center}
\begin{tabular}{@{}lrr@{}}
\toprule
 & dye-free & dot \\
\midrule
$A_{ij} = -\log_{10} R_{ij}$ & $-0.3289$ & $+0.2228$ \\
$B_{ij} = -\log_{10} S_{ij}$ & $-0.3412$ & $-0.3364$ \\
\bottomrule
\end{tabular}
\end{center}

\textbf{Why convert here and not earlier?}  Every step from here on is a
subtraction.  The post-image could be converted right after step 3.  The
pre-image has to wait until after step 4, because lining up averages light,
which is reflectance (see step 4).  So the chain converts both here, once,
and works only in absorbance after this point.

\textbf{Code today.}  \code{padphot.absorbance}, but applied today after a
division by \code{R\_white} (see Section~\ref{sec:open}).

\status{in use.}

%----------------------------------------------------------------------
\subsection{Step 6: remove the shading to reveal the stain \texorpdfstring{$\mC$}{C}}

\textbf{In words.}  The pre-image holds the pad's own shading: the shadow
at its edge and the paper's fiber pattern.  Subtracting it, in absorbance,
leaves the stain.  The method expected (method $\MD$): mix the pre-image
with a smoothed copy of itself, and subtract the mix.

\textbf{Equations.}
\begin{align}
  \widehat B_{ij} &= \lambda\, B_{ij} + (1 - \lambda)\, \widetilde B_{ij} \\
  C_{ij} &= A_{ij} - \widehat B_{ij}
\end{align}
$\tB$ is $\mB$ smoothed over width $\sigma$.  By rule 1
(Section~\ref{sec:rules}) the smoothing should average reflectance:
smooth $\mS$, then take $-\log_{10}$.  The study code so far averages
absorbance; the two differ by up to 0.008 within 4 pixels of the pad's edge
on C1032.
The intended smoothing:
two-dimensional (a bell-shaped weighted average) on the inner pad, and
along rows or columns only, at width $2\sigma$, near the pad's edges.
$\lambda$ is a \SCALAR{} between 0 and 1.
How $\lambda$ is set is open: it may be found once and stored, or set
for each photo from measures of its noise and of how well the two photos
line up (Bill, 2026-10-03).
The weights $\lambda$ and $1 - \lambda$ sum to 1,
so the same quantity can be written in reflectance with a \emph{per-pixel white}:
\begin{equation}
  \widehat S_{ij} = 10^{-\widehat B_{ij}}, \qquad T_{ij} = R_{ij} / \widehat S_{ij},
  \qquad C_{ij} = -\log_{10} T_{ij}.
\end{equation}
There is no single white number anywhere in this step.  Every pixel is
divided by its own white value $\widehat S_{ij}$, taken from the same strip
before exposure.  $C = 0$ means ``as the pad was before exposure, here''.

\textbf{What the gray-chip level does in $\mC$.}  Take $\lambda = 1$ for
clarity ($\mC = \mA - \mB$; the smoothing changes nothing below, because
$\mG$ is itself smooth).  Write $'$ for the pre-image's value at the
position the pixel came from.  From steps 3 and 5,
$A_{ij} = -\log_{10} F_{ij} + \log_{10} G_{ij}$ and
$B_{ij} \approx -\log_{10} F'_{ij} + \log_{10} G'_{ij}$, so
\begin{equation}
  C_{ij} = -\log_{10}\frac{F_{ij}}{F'_{ij}}
          + \log_{10}\frac{G_{ij}}{G'_{ij}},
  \qquad
  \frac{G_{ij}}{G'_{ij}} = \frac{P_{ij}}{P'_{ij}} \cdot
                           \frac{\gamma_j}{\gamma'_j}.
\end{equation}
\begin{itemize}
  \item \textbf{The polynomial's part}, $P_{ij}/P'_{ij}$, is the light's
    shape at two places a fraction of a pixel apart.  On C1032,
    $\log_{10}(P/P')$ is 0.00001 (range 0.00004 over the pad): it cancels.
  \item \textbf{The chips' part}, $\gamma_j/\gamma'_j$, is how much more
    light the post-image's chips collected than the pre-image's: the
    exposure time and the lamp's brightness together.  On C1032 it is 2.218
    ($\log_{10}$: $+0.346$), the same over the whole pad within 0.0001.  The
    exposure times alone give 2193/1000 = 2.193 ($+0.341$); the other 1.1\,\%
    ($+0.005$) is the lamp.
  \item So for the stain, \textbf{the gray chips act as a monitor of
    exposure and lamp between the two photos}.  Without them, every pixel's
    $\mC$ would be 0.346 too low on C1032, more than half the dot's stain.
    Neither the chips' own gray value nor the polynomial's shape reaches
    $\mC$, provided the strip moves less than a pixel.  The chips'
    left-to-right balance matters only through the ratio of the two photos:
    post/pre is 2.2188 at the left chip and 2.2172 at the right, so a single
    averaged chip would change $\mC$ by at most 0.0003 here.
  \item The shape $\mP$ does matter for anything that compares one place on
    the stage with another \emph{within one photo}.  In the chain, nothing
    after step 6 does.
\end{itemize}

\textbf{Worked numbers} (provisional settings: $\lambda = 0.81$, the value
fitted on C1032 in \file{shading\_removal.tex}; plain 1-pixel smoothing
everywhere, not yet the edge-aware smoothing).

\medskip
\begin{center}
\begin{tabular}{@{}lrr@{}}
\toprule
 & dye-free & dot \\
\midrule
$\widetilde B_{ij}$ ($\sigma = 1$) & $-0.3418$ & $-0.3399$ \\
$\widehat B_{ij} = 0.81\,B_{ij} + 0.19\,\widetilde B_{ij}$ & $-0.3413$ & $-0.3370$ \\
$C_{ij} = A_{ij} - \widehat B_{ij}$ & $+0.0124$ & $+0.5599$ \\
$\widehat S_{ij} = 10^{-\widehat B_{ij}}$ & 2.1945 & 2.1729 \\
$T_{ij} = 10^{-C_{ij}}$ & 0.9718 & 0.2755 \\
\bottomrule
\end{tabular}
\end{center}

\noindent The dye-free pixel is 2.8\,\% darker than before exposure; the
dot pixel returns about a quarter of the light it did.  Check, dot pixel:
$\widehat B = 0.81(-0.3364) + 0.19(-0.3399) = -0.3370$, and
$C = 0.2228 + 0.3370 = 0.5599$ (rounding in the last digit).

\textbf{Code today.}  Study code only: \code{shading\_review.r3\_methods}
and \code{smooth}.  Not yet in the chain.

\textbf{Scalars.}  $\lambda$, $\sigma$, $w$.  Whether one $\lambda$ suits the
whole pad is open: on C1032 the best value was 0.24 on the inner pad and
0.85 at the rim; but that used an older smoothing algorithm.

\status{under study (\file{shading\_removal.tex}).  The method, the
smoothing near the edge, and $\lambda$ are to be verified.}

%----------------------------------------------------------------------
\subsection{Step 7: the table of all channels}

\textbf{In words.}  Put every channel's stain side by side: one row per pad
pixel, one column per channel.  Only channels the pre-image holds can have
a $\mC$ (six shared by all 18 dot strips: violet.B, royal\_blue.B, cyan.B,
green.Gb, green.Gr, photo\_red.R).

\textbf{Equation.}  $\tC$ ($N \times K$), with element $C_{n,k}$ = the
stain of pad pixel $n$ in channel $k$.  Row $n$ is $\vc_n$, that pixel's
spectrum.  A region's table is written $\reg{inner}{\tC}$.

\textbf{Code today.}  \code{duotone.PadCube} (today built from
\code{R\_white}-referenced absorbance, not $\mC$).

\status{in use, to be rebuilt on $\mC$.}

%----------------------------------------------------------------------
\subsection{Step 8: yellow and brown}

\textbf{In words.}  Each pixel's spectrum is taken to be a mix of a yellow
and a brown product, plus noise.  Find how much of each.

\textbf{Equation} (the model of \file{duotone\_guide.tex}, in this
notation):
\begin{equation}
  \vc_n = h_n \vone + u_n \epsu + v_n \epsv + \text{noise},
  \qquad u_n, v_n \ge 0
\end{equation}
$\epsu$, $\epsv$: the yellow and brown spectra (lists over channels, the
same for every pixel).  $u_n$, $v_n$: the amounts at pixel $n$, painted back
as images $\mU$, $\mV$.  $h_n$: an amount that is the same in every channel.
Once step 6 has removed the shading, $h_n$ should be small; whether the
model still needs it is a question for the duotone review.
Decisions remaining:
Whether to store a single transformation matrix for use on all images (probable), or whether to recompute the transformation matrix per sample.

\textbf{Code today.}  \code{duotone.unmix}, \code{endmembers},
\code{contrast}.

\status{in code with known faults; Bill's review pending.  No worked
number until then.}

%----------------------------------------------------------------------
\subsection{Step 9: where the yellow and brown are}

\textbf{In words.}  The sizes and shapes of the yellow and brown areas may
carry more information than their amounts alone.

\status{not started.  Depends on the exposure geometry.}

%----------------------------------------------------------------------
\subsection{Step 10: relate the measures to arsenic}

\textbf{In words.}  Fit the per-strip measures from steps 8--9 against the
known concentrations $x_\ell$ of the calibration standards, then invert the
fit to estimate an unknown.

\status{not started.  The fit found with the current geometry may not carry
over to the final one.}

%======================================================================
\section{What is not the chain}
\label{sec:inventory}

Files in \file{desktop/padanalysis/}, sorted by role.  ``Chain'' means it
does a step of Section~\ref{sec:chain} today (often among other things).

\medskip
\begin{tabular}{@{}L{2.6cm}L{12.2cm}@{}}
\toprule
Role & Files and functions \\
\midrule
Chain &
  \code{probe.Capture} (\code{layer}, \code{dark\_level}, \code{hot\_mask});
  \code{analysis\_1}: \code{open\_sample}, \code{rung\_verdict},
  \code{\_region}, \code{\_hot}, \code{pad\_of}, \code{pre\_to\_post},
  \code{\_register}, \code{get\_balanced\_cell\_data} (up to $\mR$);
  \code{padphot}: \code{divisor\_map}, \code{reflectance},
  \code{absorbance};
  \code{whitecorr.LayerModel}, \code{load\_models};
  \code{padrect.PadRect} and \code{fields} (pad finding);
  \code{duotone}: \code{PadCube}, \code{unmix}, \code{endmembers},
  \code{contrast} (step 8, under review);
  \code{pca\_study.dot\_geometry}, \code{dot\_study.dot\_edge} (dot, provisional).
  In \file{src/}: \code{capture\_h5.dark\_level}, \code{hotpixel\_repo},
  \code{darkframe}. \\
\midrule
Inspection &
  \code{analysis\_1}: \code{roi\_overlay}, \code{contact\_sheet},
  \code{cell\_table}, \code{pre\_post\_overlay}, \code{inventory},
  \code{cell\_PCA}, the interim store (\code{writeh5\_sampcell} \ldots);
  \code{colorview}; \code{chain\_example}; \code{duotone.render\_rgb}. \\
\midrule
Research &
  \code{shading\_review}, \code{shading\_study}, \code{dot\_study},
  \code{pca\_study}, \code{pca}, \code{padregistration},
  \code{analysis\_2}, \code{hotpixel\_analysis}, \code{duotone\_guide},
  \code{R\_guide}, \code{R\_inf\_guide};
  \code{analysis\_1.pre\_white} (\code{R\_white}) and the $R_\infty$,
  K/S quantities in \code{padphot}
  (\code{reflectance\_infinity}, \code{km\_remission} \ldots); duotone's
  noise and channel-choice tools (\code{amount\_sigma}, \code{noise\_cov},
  \code{j\_hue}, \code{j\_amounts}, \code{subset\_*}). \\
\midrule
History &
  \code{report1}--\code{report14}, \code{measure}, \code{chemistry},
  \code{sweep}, \code{look1}--\code{look3}, \code{segment}.  Built for the
  whole-pad strips C1002--C1014, most of which have no pre-image. \\
\bottomrule
\end{tabular}

\medskip
\noindent The plan (agreed 2026-10-03) is a new, small chain module with one
function per step above, named by these symbols and reusing \code{probe},
\code{padrect} and \code{padphot}.  \code{analysis\_1} stays as the
inspection and research workbench.

%======================================================================
\section{Open points, with recommendations}
\label{sec:open}

\noindent Bill accepted these recommendations on 2026-10-03, except where a
point says otherwise.

\begin{enumerate}[leftmargin=1.5em]
  \item \textbf{The single white value \code{R\_white} is not part of the
    chain.}  Today the code divides every pixel by one number per channel:
    the median of the pre-image pad (\code{analysis\_1.pre\_white}).  In
    step 6 the white is an image, $\hS$, so that number cancels exactly and
    is not needed.  It stays only in displays and side studies that use it.
    The current workbench (\code{analysis\_1}, \code{duotone.PadCube}) keeps
    it until the new chain module replaces it.  If shading removal ever
    settled on a method that needs a single white value (method $\MA$),
    this would come back to Bill.  \emph{Recommendation:} after the rewrite,
    check together that no single white value remains in the chain, and
    look at $\hS$ on a pad (map and histogram).
  \item \textbf{Dark level, wording.}  The code subtracts one number per
    Bayer layer per photo: the mean of that photo's \emph{shortest-exposure}
    dark frame.  Dark frames are shared by exposure time, not by LED.  Your
    outline said ``a current matching dark image channel (LED + Bayer
    layer)''.  The two agree in effect, because the dark level depends on
    neither the LED nor the exposure (64.02--64.04\,DN).
    \emph{Decided:} the chain is stated as ``one dark-frame mean per Bayer
    layer per photo''.
  \item \textbf{Blending several exposures} (step 1).  Mentioned long ago,
    never examined.  \emph{Bill to review the idea first.}
  \item \textbf{Noise lowered by lining up} (step 4).  The 4-pixel average
    makes the pre-image look less noisy than it is.
    \emph{Measured 2026-10-03} (\file{shading\_removal.tex}, Section 4):
    over 18 pads and six channels the lined-up pre-image keeps a median
    0.61 of its pixel-to-pixel scatter (range 0.43--0.96, set by where each
    pad landed between pixels); lining up by nearest pixel keeps all of it.
    Noise must be measured on the pre-image's own grid (the yellow/brown
    code already does).
  \item \textbf{One $\lambda$ for the whole pad} (step 6).  C1032 wants 0.24
    inside and 0.85 at the rim.  \emph{Bill's position:} a $\lambda$ that varies
    risks being confused with the stain.  First try to make one $\lambda$
    work through the right smoothing near the edge.  Only if the outermost
    one to three rows along each pad edge are shaded in a fundamentally
    different way, let $\lambda$ vary there, and only there.
  \item \textbf{Two sources of concentrations} (step 10).
    \file{samples.py} holds $x_\ell$ for C1002--C1014; the dot strips'
    concentrations are in the database (\code{samples.manual\_read}).
    \emph{Decided:} the chain reads only the database.  Bill asks that
    every other small, special-purpose data store be called out, because
    that is not the desired practice.  Found in
    \file{desktop/padanalysis/}:
    \begin{itemize}[nosep]
      \item \file{samples.py}: a hand-written table of strips, concentrations
        and preparation flags (C1002--C1014);
      \item \code{probe.FILES}: a second copy of concentrations for
        X15 and C1002--C1005, used by old reports and \code{segment};
      \item \file{records.json}, \file{records\_blob.json},
        \file{ratios\_by\_erode.json}: results files written and read by
        the old reports (\code{measure}, \code{chemistry}, \code{report*});
      \item \code{duotone.NM}: LED peak wavelengths, copied by hand from
        \file{src/lighting\_cal8.py};
      \item \code{analysis\_1.GRAYCELLS} and \code{PRECELLS}: channel lists
        typed in from two captures (\code{cells\_present} can regenerate
        them);
      \item \file{analysis/samples/<sample>.h5}: an interim cache of
        computed arrays, stamped with the chain version.
    \end{itemize}
    The first three belong to the history pipeline.  The new chain module
    should read concentrations, wavelengths and channel lists from the
    database or the capture file itself, and keep any cache only as a
    cache.
  \item \textbf{Older reports use other letters.}
    \file{colorimeter\_flatfield.tex} uses $A$ for reflectance and $S$ for
    signal, and \file{duotone\_guide.tex} uses $y$, $b$, $g$ for the amounts.
    \emph{Recommendation:} leave settled reports as they are, with a
    one-line pointer to this notation; rewrite only reports still being
    worked on (\file{shading\_removal.tex}, then \file{duotone\_guide.tex}
    at its review).
  \item \textbf{K/S} is computed for every channel but used by no step of
    the chain.  \emph{Recommendation:} research only; leave it out of the
    chain module.
  \item \textbf{Smoothing the pre-image in reflectance} (step 6, rule 1).
    The study code smooths $\mB$ (absorbance); rule 1 says smooth $\mS$
    (reflectance) and then take $-\log_{10}$.  They agree inside the pad and
    differ by up to 0.008 within 4 pixels of the edge on C1032.
    \emph{Recommendation:} switch to smoothing $\mS$ in the edge-smoothing
    step, and compare.
\end{enumerate}

%======================================================================
\section{How to reproduce}

\begin{verbatim}
cd desktop/padanalysis
timeout -k 10 300 ../../.venv/bin/python3 -u chain_example.py
\end{verbatim}

\noindent prints every worked number in Section~\ref{sec:chain}.  It uses
the chain's current functions, removes \code{R\_white} from $\mA$ and $\mB$
(by adding back $\log_{10}$\code{R\_white}), and applies step 6 with the
provisional settings stated there.

\end{document}
